% =============================================================================
% APPENDIX DP: REF-DEVPIPE - Developer Pipeline & Computational Architecture
% =============================================================================
% Category: REF (Reference)
% Language: English
% Version: 1.0 (January 2026)
% Status: Active
%
% Purpose: Comprehensive guide for developers showing how EBF components
%          connect computationally, what constitutes a "result", and how
%          to implement the framework in practice.
%
% Target Audience: Developers, Programmers, Data Scientists, Engineers
%
% Complements:
%   - FM (REF-FORMALMAP): Where are the proofs?
%   - MS (REF-MODELSPACE): Which theory matches?
%   - G (REF-GLOSSARY): What do symbols mean?
%
% =============================================================================

\chapter{REF-DEVPIPE: Developer Pipeline \& Computational Architecture}
\label{app:ref-devpipe}

% =============================================================================
% ABSTRACT
% =============================================================================
\begin{abstract}
This appendix provides a comprehensive guide for developers implementing the
Evidence-Based Framework (EBF). It answers the fundamental questions: What is
a ``result'' in EBF? How do the components connect computationally? What is
the data flow from input to output? The appendix documents the complete
computational pipeline, database architecture, script inventory, and provides
concrete implementation workflows. While Appendix UNMAPPED_FM shows where proofs are,
this appendix shows how to \textit{compute} results.
\end{abstract}

% =============================================================================
% QUICK REFERENCE
% =============================================================================
\begin{tcolorbox}[colback=yellow!10!white, colframe=orange!75!black,
    title=\textbf{Quick Reference: Developer Pipeline}]

\textbf{What is a ``Result'' in EBF?}
\begin{itemize}[nosep]
    \item A \textbf{Prediction}: $P_{eff}$ for behavior change
    \item An \textbf{Intervention Design}: Optimized portfolio
    \item A \textbf{Model Output}: Calibrated 10C model
    \item A \textbf{Validation}: Theory-data comparison
\end{itemize}

\textbf{Core Pipeline:}
\begin{center}
\fbox{Input (10C) $\rightarrow$ Model $\rightarrow$ Calibration $\rightarrow$ Prediction $\rightarrow$ Validation}
\end{center}

\textbf{Key Scripts:}
\texttt{design-model}, \texttt{theory\_papers.py}, \texttt{theory\_learning.py}

\end{tcolorbox}

% =============================================================================
% HEADER BLOCK
% =============================================================================
\begin{tcolorbox}[colback=blue!5!white, colframe=blue!75!black,
    title=\textbf{Appendix UNMAPPED_DP: REF-DEVPIPE}]

\begin{tabular}{@{}ll@{}}
\textbf{Appendix:} & DP (Developer Pipeline) \\
\textbf{Category:} & REF (Reference) \\
\textbf{Purpose:} & Show how EBF computes results \\
\textbf{Audience:} & Developers, Programmers, Data Scientists \\
\textbf{Version:} & 1.0 \\
\textbf{Databases:} & 5 connected data stores \\
\end{tabular}

\tcblower

\textbf{Core Contribution:} Provides the computational blueprint for
implementing EBF, showing data flows, database connections, and concrete
code examples.

\textbf{Key Distinction:} FM shows \textit{where proofs are};
DP shows \textit{how to compute results}.

\end{tcolorbox}

% =============================================================================
% APPENDIX SCOPE BOX
% =============================================================================
\begin{tcolorbox}[
    colback=orange!5!white,
    colframe=orange!75!black,
    title={\textbf{Appendix Scope: Ziel / In-Scope / Out-of-Scope / Constraints}},
    fonttitle=\bfseries
]

\textbf{Ziel (Objective):}
\begin{itemize}[nosep]
    \item Define what constitutes a ``result'' in EBF
    \item Document the computational pipeline end-to-end
    \item Show how databases and scripts connect
\end{itemize}

\textbf{In-Scope:}
\begin{itemize}[nosep]
    \item 5 database architectures and their connections
    \item Complete script inventory with usage examples
    \item 4 result types and how to produce them
    \item Developer workflows with code snippets
\end{itemize}

\textbf{Out-of-Scope (delegiert an andere Appendices):}
\begin{itemize}[nosep]
    \item Mathematical proofs $\rightarrow$ Appendix UNMAPPED_FM (REF-FORMALMAP)
    \item Theory catalog $\rightarrow$ Appendix UNMAPPED_MS (REF-MODELSPACE)
    \item Parameter values $\rightarrow$ Appendix UNMAPPED_BBB (CORE-WHERE)
\end{itemize}

\textbf{Lieferobjekte (Deliverables):}
\begin{itemize}[nosep]
    \item Database architecture diagram
    \item Script reference table
    \item 4 implementation workflows
    \item Code examples for each result type
\end{itemize}

\end{tcolorbox}

% =============================================================================
% CROSS-REFERENCE MAP
% =============================================================================
\begin{tcolorbox}[colback=gray!5!white, colframe=gray!75!black,
    title=\textbf{Cross-Reference Map}]

\textbf{Outgoing References:}
\begin{itemize}[nosep]
    \item \textbf{FM (REF-FORMALMAP):} Proof locations
    \item \textbf{MS (REF-MODELSPACE):} Theory matching
    \item \textbf{BBB (CORE-WHERE):} Parameter values
    \item \textbf{EEE (METHOD-DESIGN):} Model design workflow
\end{itemize}

\textbf{Incoming References:}
\begin{itemize}[nosep]
    \item All developers implementing EBF
    \item Data scientists building models
    \item Engineers creating production systems
\end{itemize}

\end{tcolorbox}

% =============================================================================
% CHAPTER LINKAGE
% =============================================================================
\begin{tcolorbox}[colback=purple!5!white, colframe=purple!75!black,
    title=\textbf{Chapter Linkage}]

\textbf{Primary Chapter:} Chapter 8 (Mathematical Foundations)

\textbf{Secondary Chapters:}
\begin{itemize}[nosep]
    \item Chapter 16: Probability of Behavior Change (prediction output)
    \item Chapter 17: Intervention Design (intervention output)
\end{itemize}

\end{tcolorbox}

% =============================================================================
% MOTIVATION
% =============================================================================
\section{Motivation: Why Do We Need This Appendix?}
\label{sec:dp-motivation}

\begin{tcolorbox}[colback=orange!5!white, colframe=orange!75!black,
    title=\textbf{Motivation: The Developer's Dilemma}]

\textbf{The Problem:}

A developer joins the EBF project and asks:
\begin{quote}
``I understand the theory, but how do I actually \textit{compute} something?
What databases exist? What scripts should I use? What is even a `result'?''
\end{quote}

\textbf{The Gap:}

\begin{center}
\begin{tabular}{|l|l|l|}
\hline
\textbf{Appendix} & \textbf{Question Answered} & \textbf{Missing} \\
\hline
\hline
FM & Where is the proof? & --- \\
MS & Which theory matches? & --- \\
G & What does symbol mean? & --- \\
\textbf{???} & \textbf{How do I compute?} & \textbf{DP fills this gap} \\
\hline
\end{tabular}
\end{center}

\textbf{What DP Provides:}
\begin{itemize}[nosep]
    \item Definition of ``result'' types
    \item Database architecture and connections
    \item Script inventory with examples
    \item End-to-end implementation workflows
\end{itemize}

\end{tcolorbox}

% =============================================================================
% FUNDAMENTAL QUESTION
% =============================================================================
\section{The Fundamental Question}
\label{sec:dp-fundamental}

\begin{tcolorbox}[colback=red!5!white, colframe=red!75!black,
    title=\textbf{The Fundamental Question}]

\textbf{Question:} What is a ``result'' in EBF and how do I compute it?

\textbf{Answer:} A result is a \textbf{validated output} from the computational pipeline:
\begin{equation}
\text{Input (10C Model)} \xrightarrow{\text{calibrate}} \text{Parameters} \xrightarrow{\text{compute}} \text{Output} \xrightarrow{\text{validate}} \textbf{Result}
\end{equation}

\end{tcolorbox}

% =============================================================================
% CORE THEORY: COMPUTATIONAL SEMANTICS
% =============================================================================
\section{Core Theory: Computational Semantics of EBF}
\label{sec:dp-theory}

\begin{tcolorbox}[colback=blue!3!white, colframe=blue!60!black,
    title=\textbf{Theoretical Foundation: What Makes EBF Computable?}]

\textbf{Theorem UNMAPPED_DP-1 (Computability):}
Every EBF model with fully specified 10C dimensions yields a computable prediction.

\textbf{Proof sketch:}
\begin{enumerate}[nosep]
    \item 10C dimensions are finite and enumerable
    \item Each dimension maps to bounded parameters ($\Theta \in \mathbb{R}^n$)
    \item Utility functions are continuous and bounded
    \item $P_{eff}$ is a well-defined function of 10C $\Rightarrow$ computable
\end{enumerate}

\textbf{Theorem UNMAPPED_DP-2 (Validation Convergence):}
With sufficient project results, theory parameters converge to true values.

\textbf{Formal statement:}
\begin{equation}
\lim_{n \rightarrow \infty} \hat{\Theta}_n = \Theta^* \quad \text{(in probability)}
\end{equation}
where $\hat{\Theta}_n$ is the estimate after $n$ project results.

\textbf{Theorem UNMAPPED_DP-3 (Pipeline Consistency):}
The 6-stage pipeline preserves information:
\begin{equation}
\text{Input}(I) \xrightarrow{\text{Pipeline}} \text{Result}(R) \Rightarrow I \text{ recoverable from } R
\end{equation}

\end{tcolorbox}

% =============================================================================
% SECTION 1: WHAT IS A RESULT?
% =============================================================================
\section{What is a ``Result'' in EBF?}
\label{sec:dp-results}

\begin{tcolorbox}[colback=blue!5!white, colframe=blue!75!black,
    title=\textbf{The Four Result Types}]

\textbf{Type 1: PREDICTION}
\begin{itemize}[nosep]
    \item \textbf{Output:} $P_{eff}$ (probability of behavior change)
    \item \textbf{Formula:} $P_{eff} = f(\text{AWARE}, \text{READY}, \text{STAGE}, \text{Segment}, \Psi)$
    \item \textbf{Script:} \texttt{/r-score --demo}
    \item \textbf{Validation:} Compare to observed behavior rates
\end{itemize}

\vspace{0.5em}
\textbf{Type 2: INTERVENTION DESIGN}
\begin{itemize}[nosep]
    \item \textbf{Output:} Optimized intervention portfolio
    \item \textbf{Formula:} $\text{Portfolio}^* = \arg\min \sum_i \text{cost}_i$ s.t. $P \geq P_{target}$
    \item \textbf{Script:} \texttt{/design-intervention}
    \item \textbf{Validation:} A/B testing, pilot results
\end{itemize}

\vspace{0.5em}
\textbf{Type 3: MODEL OUTPUT}
\begin{itemize}[nosep]
    \item \textbf{Output:} Calibrated 10C model with parameters $\Theta$
    \item \textbf{Formula:} All 10C dimensions specified with values
    \item \textbf{Script:} \texttt{/design-model}
    \item \textbf{Validation:} Theory match (MS), parameter bounds (BBB)
\end{itemize}

\vspace{0.5em}
\textbf{Type 4: VALIDATION}
\begin{itemize}[nosep]
    \item \textbf{Output:} Theory-data comparison with deviation analysis
    \item \textbf{Formula:} $\Delta = |\text{predicted} - \text{observed}|$
    \item \textbf{Script:} \texttt{theory\_learning.py --project-result}
    \item \textbf{Validation:} Updates theory catalog if $\Delta > $ threshold
\end{itemize}

\end{tcolorbox}

\begin{tcolorbox}[colback=green!5!white, colframe=green!75!black,
    title=\textbf{Result Quality Criteria}]

A result is \textbf{valid} if it satisfies:

\begin{center}
\begin{tabular}{|l|l|l|}
\hline
\textbf{Criterion} & \textbf{Check} & \textbf{Tool} \\
\hline
\hline
Theory Grounded & Matches established theory & MS lookup \\
Parameter Bounded & Within literature ranges & BBB check \\
Internally Consistent & No contradictions in 10C & Model validation \\
Externally Validated & Compared to data & Project results \\
\hline
\end{tabular}
\end{center}

\end{tcolorbox}

% =============================================================================
% SECTION 2: DATABASE ARCHITECTURE
% =============================================================================
\section{Database Architecture}
\label{sec:dp-databases}

\begin{tcolorbox}[colback=blue!3!white, colframe=blue!60!black,
    title=\textbf{The 5 Connected Databases}]

\begin{center}
\begin{tabular}{|l|l|l|l|}
\hline
\textbf{\#} & \textbf{Database} & \textbf{File} & \textbf{Content} \\
\hline
\hline
1 & \textbf{Paper Database} & \texttt{bcm\_master.bib} & 1,922+ papers \\
2 & \textbf{Theory Catalog} & \texttt{theory-catalog.yaml} & 134 models \\
3 & \textbf{Case Registry} & \texttt{case-registry.yaml} & Historical cases \\
4 & \textbf{Intervention Registry} & \texttt{intervention-registry.yaml} & Projects \\
5 & \textbf{Concept Registry} & \texttt{concept-registry.yaml} & EIP tracking \\
\hline
\end{tabular}
\end{center}

\end{tcolorbox}

\begin{tcolorbox}[colback=purple!3!white, colframe=purple!60!black,
    title=\textbf{Database Connection Architecture}]

\begin{verbatim}
                    ┌─────────────────────┐
                    │   bcm_master.bib    │
                    │   (1,922+ papers)   │
                    └──────────┬──────────┘
                               │ theory_support
                               ▼
                    ┌─────────────────────┐
                    │  theory-catalog.yaml │◄────────────────┐
                    │   (134 models)      │                  │
                    └──────────┬──────────┘                  │
                               │ bib_keys                    │
            ┌──────────────────┼──────────────────┐          │
            ▼                  ▼                  ▼          │
┌─────────────────┐ ┌─────────────────┐ ┌─────────────────┐  │
│ case-registry   │ │ intervention-   │ │ concept-registry│  │
│    .yaml        │ │ registry.yaml   │ │    .yaml        │  │
│ (historical)    │ │ (projects)      │ │ (EIP tracking)  │──┘
└─────────────────┘ └─────────────────┘ └─────────────────┘
        │                   │                   │
        └───────────────────┴───────────────────┘
                            │
                    ┌───────▼───────┐
                    │  LEARNING LOOP │
                    │ (auto-updates) │
                    └───────────────┘
\end{verbatim}

\textbf{Data Flows:}
\begin{enumerate}[nosep]
    \item \textbf{Paper $\rightarrow$ Theory:} \texttt{theory\_support} field links papers to theories
    \item \textbf{Theory $\rightarrow$ Paper:} \texttt{bib\_keys} field links theories to papers
    \item \textbf{Project $\rightarrow$ Theory:} Results update theory parameters
    \item \textbf{Concept $\rightarrow$ Theory:} New concepts add to catalog via EIP
\end{enumerate}

\end{tcolorbox}

% =============================================================================
% SECTION 3: SCRIPT INVENTORY
% =============================================================================
\section{Script Inventory}
\label{sec:dp-scripts}

\begin{tcolorbox}[colback=green!3!white, colframe=green!60!black,
    title=\textbf{Complete Script Reference}]

\small
\begin{longtable}{|l|l|l|}
\hline
\textbf{Script} & \textbf{Purpose} & \textbf{Example} \\
\hline
\hline
\endfirsthead
\hline
\textbf{Script} & \textbf{Purpose} & \textbf{Example} \\
\hline
\endhead
\multicolumn{3}{|c|}{\textbf{Model Design}} \\
\hline
\texttt{/design-model} & Create 10C model & \texttt{/design-model --mode schnell} \\
\texttt{/design-intervention} & Design intervention & \texttt{/design-intervention} \\
\hline
\multicolumn{3}{|c|}{\textbf{Calibration (LLMMC)}} \\
\hline
\texttt{/r-score} & LLM Monte Carlo estimation & \texttt{/r-score --demo} \\
\texttt{/r-score --full} & Full LLMMC pipeline & \texttt{/r-score --full} \\
\hline
\multicolumn{3}{|c|}{\textbf{Theory \& Paper Lookup}} \\
\hline
\texttt{theory\_papers.py} & Theory $\leftrightarrow$ Paper & \texttt{--theory MS-RD-001} \\
\texttt{search\_bibliography.py} & Search papers & \texttt{--author "Thaler"} \\
\texttt{assign\_lit\_appendix.py} & Assign LIT category & \texttt{-a "Grant" -t "..."} \\
\hline
\multicolumn{3}{|c|}{\textbf{Output Generation (8D)}} \\
\hline
\texttt{generate\_paper.py} & Generate document & \texttt{--style fehr} \\
\texttt{/board-presentation} & Board slides & \texttt{Company pdf} \\
\hline
\multicolumn{3}{|c|}{\textbf{Learning Loop}} \\
\hline
\texttt{theory\_learning.py} & Update theories & \texttt{--project-result} \\
\texttt{sync\_bib\_to\_lit.py} & Sync BibTeX to LIT & \texttt{--update} \\
\hline
\multicolumn{3}{|c|}{\textbf{Validation}} \\
\hline
\texttt{check\_template\_compliance.py} & Appendix compliance & \texttt{appendices/*.tex} \\
\texttt{check\_chapter\_compliance.py} & Chapter compliance & \texttt{chapters/*.tex} \\
\texttt{validate\_core\_framework.py} & 10C validation & (no args) \\
\hline
\multicolumn{3}{|c|}{\textbf{Case \& Intervention}} \\
\hline
\texttt{/case} & Query cases & \texttt{/case --domain health} \\
\texttt{/intervention} & Query interventions & \texttt{/intervention --learnings} \\
\hline
\end{longtable}

\end{tcolorbox}

% =============================================================================
% SECTION 4: COMPUTATIONAL PIPELINE
% =============================================================================
\section{The Computational Pipeline}
\label{sec:dp-pipeline}

\begin{tcolorbox}[colback=blue!5!white, colframe=blue!75!black,
    title=\textbf{End-to-End Pipeline: From Question to Result}]

\begin{verbatim}
┌─────────────────────────────────────────────────────────────────────┐
│ STAGE 1: PROBLEM DEFINITION                                         │
│ ─────────────────────────────────────────────────────────────────── │
│ Input: Business question ("How to increase savings rate?")          │
│ Output: 10C specification (WHO, WHAT, HOW, WHEN, WHERE, ...)        │
│ Script: /design-model --mode geführt                                │
└─────────────────────────────────────────────────────────────────────┘
                                    │
                                    ▼
┌─────────────────────────────────────────────────────────────────────┐
│ STAGE 2: THEORY MATCHING                                            │
│ ─────────────────────────────────────────────────────────────────── │
│ Input: 10C specification                                            │
│ Output: Matched theories (e.g., MS-TD-002 Hyperbolic Discounting)   │
│ Script: python theory_papers.py --match-10c "beta < 1"              │
└─────────────────────────────────────────────────────────────────────┘
                                    │
                                    ▼
┌─────────────────────────────────────────────────────────────────────┐
│ STAGE 3: PARAMETER CALIBRATION                                      │
│ ─────────────────────────────────────────────────────────────────── │
│ Input: Matched theories                                             │
│ Output: Parameter values Θ                                          │
│ Methods:                                                            │
│   (a) BBB Literature: python theory_papers.py --theory MS-TD-002 -v │
│   (b) LLMMC: /r-score --demo (LLM Monte Carlo estimation)           │
│   (c) Empirical: Custom data calibration                            │
│ Check: Parameters within literature bounds (BBB)                    │
└─────────────────────────────────────────────────────────────────────┘
                                    │
                                    ▼
┌─────────────────────────────────────────────────────────────────────┐
│ STAGE 4: PREDICTION / INTERVENTION DESIGN                           │
│ ─────────────────────────────────────────────────────────────────── │
│ Input: Calibrated model                                             │
│ Output: P_eff prediction OR intervention portfolio                  │
│ Script: /r-score OR /design-intervention                            │
│ Reference: Ch 15-16 (Prediction), Ch 17-20 (Interventions)          │
└─────────────────────────────────────────────────────────────────────┘
                                    │
                                    ▼
┌─────────────────────────────────────────────────────────────────────┐
│ STAGE 5: OUTPUT GENERATION (8D Document Framework)                  │
│ ─────────────────────────────────────────────────────────────────── │
│ Input: Results from Stage 4                                         │
│ Output: Audience-appropriate document (Paper/Brief/Report/Slides)   │
│ Method: 8D Coordinates → Structure → Style → Vocabulary             │
│ Script: python scripts/generate_paper.py --style fehr               │
│ Reference: Appendix DDD (Theory), Appendix CCC (Computation)        │
│ Chapters: Ch 17 (§ Communication), Ch 22 (Policy Applications)      │
└─────────────────────────────────────────────────────────────────────┘
                                    │
                                    ▼
┌─────────────────────────────────────────────────────────────────────┐
│ STAGE 6: VALIDATION & LEARNING                                      │
│ ─────────────────────────────────────────────────────────────────── │
│ Input: Prediction + Observed outcome                                │
│ Output: Deviation analysis, theory updates                          │
│ Script: python theory_learning.py --project-result                  │
│ Loop: Results feed back to theory catalog                           │
│ Reference: Ch 21-24 (Dynamic Evolution, Policy, Multi-Level, Life)  │
└─────────────────────────────────────────────────────────────────────┘
\end{verbatim}

\end{tcolorbox}

% =============================================================================
% SECTION 4b: CALIBRATION METHODS
% =============================================================================
\subsection{The Three Calibration Methods}
\label{sec:dp-calibration}

\textbf{Navigation Hub:} For complete parameterization methodology including decision trees and parameter-type routing, see \textbf{Appendix UNMAPPED_PM (REF-PARAMETH)}.

\begin{tcolorbox}[colback=yellow!5!white, colframe=yellow!75!black,
    title=\textbf{Stage 3 Deep Dive: Parameter Calibration Methods}]

\textbf{Method A: Literature-Based (BBB)}
\begin{itemize}[nosep]
    \item \textbf{Source:} Appendix UNMAPPED_BBB (CORE-WHERE) parameter repository
    \item \textbf{Script:} \texttt{python theory\_papers.py --theory MS-TD-002 -v}
    \item \textbf{Output:} Literature-validated parameter ranges (e.g., $\beta \in [0.7, 0.95]$)
    \item \textbf{When to use:} Standard cases with established literature
\end{itemize}

\vspace{0.5em}
\textbf{Method B: LLM Monte Carlo (LLMMC)}
\begin{itemize}[nosep]
    \item \textbf{Source:} LLM-based expert simulation
    \item \textbf{Script:} \texttt{/r-score --demo} or \texttt{/r-score --full}
    \item \textbf{Output:} R-Score with confidence intervals
    \item \textbf{When to use:} Novel contexts, limited empirical data, rapid estimation
    \item \textbf{Reference:} Appendix UNMAPPED_CAL (METHOD-LLMMC)
\end{itemize}

\vspace{0.5em}
\textbf{Method C: Empirical Calibration}
\begin{itemize}[nosep]
    \item \textbf{Source:} Project-specific data
    \item \textbf{Script:} Custom analysis scripts
    \item \textbf{Output:} Context-specific parameters
    \item \textbf{When to use:} Sufficient local data available, high-stakes decisions
\end{itemize}

\vspace{0.5em}
\textbf{Decision Matrix:}

\begin{center}
\begin{tabular}{|l|c|c|c|}
\hline
\textbf{Situation} & \textbf{BBB} & \textbf{LLMMC} & \textbf{Empirical} \\
\hline
\hline
Standard domain & \checkmark & --- & --- \\
Novel context & --- & \checkmark & --- \\
High-stakes & \checkmark & \checkmark & \checkmark \\
Rapid prototype & --- & \checkmark & --- \\
Rich local data & --- & --- & \checkmark \\
\hline
\end{tabular}
\end{center}

\end{tcolorbox}

% =============================================================================
% SECTION 4c: OUTPUT GENERATION (8D)
% =============================================================================
\subsection{Stage 5: Output Generation (8D Framework)}
\label{sec:dp-output}

\textbf{Reference Appendices:} Appendix DDD (REF-DOCTYPE) for theory, Appendix CCC (METHOD-DOCTYPE) for computation.

\begin{tcolorbox}[colback=purple!5!white, colframe=purple!75!black,
    title=\textbf{Stage 5 Deep Dive: 8D Document Framework}]

\textbf{Core Principle:} Output format emerges from audience characteristics.

\vspace{0.5em}
\textbf{The 8 Dimensions:}

\begin{center}
\small
\begin{tabular}{|c|l|l|l|}
\hline
\textbf{D} & \textbf{Dimension} & \textbf{Question} & \textbf{Scale} \\
\hline
\hline
D$_1$ & Knowledge & How expert? & 0=Layperson $\rightarrow$ 1=Expert \\
D$_2$ & Proximity & How close to field? & 0=Distant $\rightarrow$ 1=Same field \\
D$_3$ & Scope & Personal or societal? & 0=Personal $\rightarrow$ 1=Societal \\
D$_4$ & Time & Reading time available? & 0=Little $\rightarrow$ 1=Much \\
D$_5$ & Goal & Communication goal? & G$_1$-G$_7$ (see below) \\
D$_6$ & Context & Internal or external? & 0=Internal $\rightarrow$ 1=Public \\
D$_7$ & Emotion & Factual or emotional? & 0=Factual $\rightarrow$ 1=Emotional \\
D$_8$ & Persistence & Ephemeral or archived? & 0=Ephemeral $\rightarrow$ 1=Archive \\
\hline
\end{tabular}
\end{center}

\vspace{0.5em}
\textbf{D$_5$ Goal Categories:}
G$_1$=Inform, G$_2$=Act, G$_3$=Convince, G$_4$=Entertain, G$_5$=Experience, G$_6$=Connect, G$_7$=Express

\vspace{0.5em}
\textbf{Output Emergence Rules:}

\begin{center}
\small
\begin{tabular}{|l|l|}
\hline
\textbf{8D Profile} & \textbf{Emergent Output} \\
\hline
\hline
D$_1$>0.7, D$_8$>0.8, D$_5$=G$_1$ & \textbf{Journal Paper} (LaTeX, hedging, citations) \\
D$_1$<0.5, D$_4$<0.3, D$_5$=G$_2$ & \textbf{Policy Brief} (Executive Summary, CTA) \\
D$_6$<0.3, D$_4$<0.2, D$_5$=G$_1$ & \textbf{Internal Memo} (Bullet points, plain) \\
D$_1$>0.6, D$_4$<0.4, D$_5$=G$_3$ & \textbf{Board Slides} (Key metrics, visuals) \\
\hline
\end{tabular}
\end{center}

\vspace{0.5em}
\textbf{Script Usage:}
\begin{verbatim}
# Generate paper in Fehr style
python scripts/generate_paper.py appendices/input.tex --style fehr

# Available styles: fehr, thaler, kahneman, sunstein
\end{verbatim}

\vspace{0.5em}
\textbf{Chapter References:}
\begin{itemize}[nosep]
    \item Ch 17 (Intervention Foundations): Communication of intervention design
    \item Ch 22 (Policy Applications): Policy-oriented output formatting
\end{itemize}

\end{tcolorbox}

% =============================================================================
% SECTION 5: WORKED EXAMPLE
% =============================================================================
\section{Worked Example: Computing a Prediction}
\label{sec:dp-example}

\begin{tcolorbox}[colback=green!5!white, colframe=green!75!black,
    title=\textbf{Worked Example: Retirement Savings Intervention}]

\textbf{Business Question:} ``How likely are employees to increase retirement savings
if we implement automatic enrollment?''

\textbf{Stage 1: Define 10C Model}
\begin{verbatim}
/design-model --mode schnell
> WHO: Employees (L1)
> WHAT: Financial utility (FEPSDE: F)
> HOW: γ(present, future) < 0 (present bias)
\end{verbatim}

\textbf{Stage 2: Match Theory}
\begin{verbatim}
python scripts/theory_papers.py --match-10c "beta < 1, psi_default"
> MS-TD-002: Quasi-Hyperbolic Discounting
> MS-NU-002: Default Effects
\end{verbatim}

\textbf{Stage 3: Get Parameters}
\begin{verbatim}
python scripts/theory_papers.py --theory MS-TD-002 -v
> β = 0.7 (Laibson 1997)
> Papers: laibson1997golden, PAP-odonoghue1999doingdoing
\end{verbatim}

\textbf{Stage 4: Compute Prediction}
\begin{verbatim}
# Using defaults with β = 0.7
P_eff(no_default) = 0.35
P_eff(with_default) = 0.85
Δ = +50 percentage points
\end{verbatim}

\textbf{Stage 5: Generate Output (8D)}
\begin{verbatim}
# Audience: HR leadership (D1=0.4, D4=0.3, D5=G2, D6=0.2)
# Emergent format: Internal memo with action items
python scripts/generate_paper.py results.tex --style memo
> Output: recommendation_memo.pdf
> Structure: Executive summary + 3 bullet points + CTA
\end{verbatim}

\textbf{Stage 6: Validate}
\begin{verbatim}
# After implementation
python scripts/theory_learning.py --project-result \
  --theory MS-NU-002 \
  --predicted 0.85 \
  --observed 0.82 \
  --deviation 3.5%
> Status: CONFIRMED (< 10% deviation)
\end{verbatim}

\textbf{Result:} Validated prediction with 3.5\% deviation feeds back to theory catalog.

\end{tcolorbox}

% =============================================================================
% SECTION 6: KEY RESULTS
% =============================================================================
\section{Key Results}
\label{sec:dp-results-section}

\begin{tcolorbox}[colback=green!5!white, colframe=green!75!black,
    title=\textbf{Key Results of This Appendix}]

\textbf{Result 1: Result Types Defined}
\begin{quote}
EBF produces four types of results: Predictions ($P_{eff}$), Intervention Designs,
Model Outputs, and Validations.
\end{quote}

\textbf{Result 2: Database Architecture Documented}
\begin{quote}
Five databases connect bidirectionally: Papers $\leftrightarrow$ Theories $\leftrightarrow$
Cases/Interventions/Concepts, with automatic learning loop.
\end{quote}

\textbf{Result 3: Pipeline Formalized}
\begin{quote}
Six-stage pipeline: Problem Definition $\rightarrow$ Theory Matching $\rightarrow$
Calibration $\rightarrow$ Prediction/Intervention $\rightarrow$ Output Generation $\rightarrow$ Validation.
\end{quote}

\textbf{Result 4: Script Inventory Complete}
\begin{quote}
All scripts documented with purpose and usage examples for developers.
\end{quote}

\end{tcolorbox}

% =============================================================================
% SECTION 6b: COMPLETE CHAPTER-PIPELINE MAPPING
% =============================================================================
\subsection{Complete Chapter-Pipeline Mapping}
\label{sec:dp-chapter-mapping}

\begin{tcolorbox}[colback=blue!3!white, colframe=blue!60!black,
    title=\textbf{All 24 Chapters Mapped to 6-Stage Pipeline}]

The 6-stage pipeline covers \textbf{all 24 chapters} of the EBF book:

\begin{center}
\small
\begin{longtable}{|c|p{4.5cm}|p{4cm}|p{3cm}|}
\hline
\textbf{Stage} & \textbf{Description} & \textbf{Chapters} & \textbf{Key Appendices} \\
\hline
\hline
\endfirsthead
\hline
\textbf{Stage} & \textbf{Description} & \textbf{Chapters} & \textbf{Key Appendices} \\
\hline
\endhead
\textbf{1} & \textbf{Problem Definition} \newline 10C Specification & Ch 5 (Complementarity) \newline Ch 9 (Context) \newline Ch 10 (FEPSDE) \newline Ch 11 (Awareness) \newline Ch 12 (Willingness) \newline Ch 13 (BCJ) \newline Ch 14 (Segments) & AAA-HI (10 CORE) \\
\hline
\textbf{2} & \textbf{Theory Matching} \newline MS Lookup & Ch 3 (Limits) \newline Ch 6 (Reference) \newline Ch 7 (Fit/Non-Concavity) & MS (Theory Catalog) \\
\hline
\textbf{3} & \textbf{Parameter Calibration} \newline PM Navigation & Ch 4 (Empirical) \newline Ch 4x (Calibration) & PM, BBB, CAL \\
\hline
\textbf{4} & \textbf{Prediction / Intervention} \newline $P_{eff}$ or Portfolio & Ch 15 (WEC Synthesis) \newline Ch 16 (Probability) \newline Ch 17 (Intervention) \newline Ch 18 (Journey-Int.) \newline Ch 19 (Segment-Int.) \newline Ch 20 (Portfolios) & S, BC, HHH, QQ \\
\hline
\textbf{5} & \textbf{Output Generation} \newline 8D Framework & Ch 17 (§Communication) \newline Ch 22 (Policy) & DDD, CCC \\
\hline
\textbf{6} & \textbf{Validation \& Learning} \newline Feedback Loop & Ch 21 (Dynamic) \newline Ch 22 (Policy) \newline Ch 23 (Multi-Level) \newline Ch 24 (Life Journeys) & Case Registry \newline Intervention Registry \\
\hline
\end{longtable}
\end{center}

\vspace{0.5em}
\textbf{Meta-Chapters (Framework Overview):}
\begin{itemize}[nosep]
    \item \textbf{Ch 1 (Introduction):} Framework motivation and overview
    \item \textbf{Ch 2 (Rationality as Stability):} Philosophical foundation
    \item \textbf{Ch 8 (Mathematical Foundations):} Cross-cutting formal basis
\end{itemize}

\vspace{0.5em}
\textbf{Coverage Statistics:}
\begin{center}
\begin{tabular}{|l|c|}
\hline
\textbf{Metric} & \textbf{Value} \\
\hline
\hline
Chapters in Pipeline & 22 / 24 (92\%) \\
Meta-Chapters & 2 (Ch 1, 2) \\
Cross-Cutting & 1 (Ch 8) \\
Total Coverage & \textbf{100\%} \\
\hline
\end{tabular}
\end{center}

\end{tcolorbox}

% =============================================================================
% SECTION 7: INTEGRATION WITH OTHER APPENDICES
% =============================================================================
\section{Integration with Other Appendices}
\label{sec:dp-integration}

\begin{tcolorbox}[colback=purple!3!white, colframe=purple!60!black,
    title=\textbf{DP in the REF Appendix Family}]

\begin{center}
\begin{tabular}{|l|l|l|}
\hline
\textbf{Appendix} & \textbf{Question} & \textbf{Developer Use} \\
\hline
\hline
\textbf{DP (This)} & How to compute? & Implementation guide \\
FM & Where is proof? & Verify formulas \\
MS & Which theory? & Match 10C to literature \\
G & What means symbol? & Understand code variables \\
BBB & What parameter value? & Calibrate models \\
\hline
\end{tabular}
\end{center}

\textbf{Typical Developer Workflow:}
\begin{enumerate}[nosep]
    \item \textbf{DP:} Understand pipeline and result types
    \item \textbf{MS:} Find matching theory for 10C model
    \item \textbf{FM:} Verify mathematical foundation
    \item \textbf{BBB:} Get parameter bounds
    \item \textbf{DP:} Execute pipeline, validate results
\end{enumerate}

\end{tcolorbox}

% =============================================================================
% SUMMARY
% =============================================================================
\section{Summary}
\label{sec:dp-summary}

\begin{tcolorbox}[colback=blue!5!white, colframe=blue!75!black,
    title=\textbf{Appendix UNMAPPED_DP Summary}]

\textbf{Core Contribution:}
\begin{itemize}[nosep]
    \item Defined 4 result types (Prediction, Intervention, Model, Validation)
    \item Documented 5-database architecture with connections
    \item Formalized 5-stage computational pipeline
    \item Provided complete script inventory
\end{itemize}

\textbf{Key Insight:}
\begin{quote}
A ``result'' in EBF is not just an output---it is a \textit{validated} output
that feeds back into the learning loop, improving future predictions.
\end{quote}

\textbf{Usage:} Developers start here to understand how to implement EBF.

\end{tcolorbox}

% =============================================================================
% GLOSSARY
% =============================================================================
\section{Glossary of Key Terms}
\label{sec:dp-glossary}

\begin{itemize}[nosep]
    \item \textbf{Result:} Validated output from computational pipeline
    \item \textbf{Pipeline:} 5-stage process from question to validated result
    \item \textbf{Learning Loop:} Automatic update of theories from project results
    \item \textbf{Calibration:} Setting parameter values from literature or data
\end{itemize}

\medskip
\textbf{Complete glossary:} See Appendix G (REF-GLOSSARY) for all EBF terminology.

% =============================================================================
% OPEN ISSUES
% =============================================================================
\section{Open Issues and Research Agenda}
\label{sec:dp-open-issues}

\begin{tcolorbox}[colback=red!5!white, colframe=red!75!black,
    title=\textbf{Open Issues}]

\textbf{Issue 1: Automation Level}
\begin{itemize}[nosep]
    \item Current: Semi-automated (human in loop)
    \item Future: Fully automated prediction pipelines
\end{itemize}

\textbf{Issue 2: API Development}
\begin{itemize}[nosep]
    \item Current: CLI scripts only
    \item Future: REST API for external integration
\end{itemize}

\textbf{Issue 3: Real-time Validation}
\begin{itemize}[nosep]
    \item Current: Post-hoc validation
    \item Future: Streaming validation with live data
\end{itemize}

\end{tcolorbox}

\textbf{Research Agenda:}
\begin{enumerate}[nosep]
    \item Develop REST API for EBF predictions
    \item Implement streaming validation
    \item Create visual pipeline editor
    \item Build automated testing suite
\end{enumerate}

% =============================================================================
% REFERENCES
% =============================================================================
\section{References}
\label{sec:dp-references}

This appendix references the master bibliography containing 1,922+ papers.

\nocite{bcm_master}

\textbf{Key implementation references:}
\begin{itemize}[nosep]
    \item All scripts in \texttt{scripts/} directory
    \item Database schemas in \texttt{data/} directory
    \item Workflow documentation in \texttt{.claude/commands/}
\end{itemize}

% =============================================================================
% CROSS-REFERENCE FOOTER
% =============================================================================
\vfill
\begin{tcolorbox}[colback=gray!10!white, colframe=gray!50!black]
\small
\textbf{Cross-References:}
FM (Proofs) $\bullet$
MS (Theories) $\bullet$
G (Glossary) $\bullet$
BBB (Parameters) $\bullet$
EEE (Model Design) $\bullet$
DR (Data Requirements)

\medskip
\textbf{Navigation:} Question $\xrightarrow{\text{10C}}$ Model $\xrightarrow{\text{calibrate}}$ Prediction $\xrightarrow{\text{validate}}$ \textbf{Result}
\end{tcolorbox}

\end{chapter}
